\section{Розділ~\ref{sec:gaba}. \nt{GLUT} і \nt{GABA} разом}

Логічні й динамічні твердження розділу виведено в ньому самому з лінійного аналізу й закону Ома; дослід показує, що схеми, на які вони спираються, --- заборона із затримкою, пряме гальмування слідом за збудженням, стабілізація гальмуванням, ритм петлі \nt{GLUT}--\nt{GABA}, --- у мозку справді працюють.

{\footnotesize
\setlength{\tabcolsep}{3pt}\renewcommand{\arraystretch}{1.15}
\begin{longtable}{>{\raggedright}p{0.5cm}>{\raggedright}p{4.0cm}>{\raggedright}p{5.6cm}>{\raggedright}p{2.3cm}>{\raggedright\arraybackslash}p{1.7cm}}
\toprule
№ & Твердження & Дослід і що виміряно & Джерело & Статус \\
\midrule\endfirsthead
\toprule
№ & Твердження & Дослід і що виміряно & Джерело & Статус \\
\midrule\endhead
1 & \nt{GABA}-нейронів у корі від п'ятої частини до чверті & Імунофарбування на GABA у $10$ областях кори мавпи: близько $25\,\%$ нейронів у $8$ областях, менше $20\,\%$ у первинній зоровій корі. Огляд різноманіття гальмівних нейронів кори & \cite{Hendry1987,Markram2004} & виміряно \\
2 & Що робить гальмування, багато в чому вирішує місце посадки: дендрит, тіло, місце народження імпульсу, закінчення чужого проводу & Огляд: класи \nt{GABA}-нейронів кори різняться мішенню на отримувачі. Одночасний запис із тіла й дендрита пірамідного нейрона: пряме гальмування набагато сильніше на тілі, ніж на дендритах. Гальмування на закінченні чужого проводу в спинному мозку зменшує викид медіатора однієї вхідної лінії (огляд вимірювань) & \cite{Markram2004,Pouille2001,Rudomin1999} & виміряно \\
3 & Зміна знака через посередника коштує однієї затримки, близько мілісекунди; гальмування, що прийшло слідом за збудженням, звужує вікно збігу & Пірамідні нейрони гіпокампа щура у зрізах: гальмування через одного посередника приходить слідом за прямим збудженням із короткою затримкою, і вікно, у якому два збуджувальні входи підсумовуються до порога, менше за $2\ms$. На дендритах, де це гальмування слабше, вікно ширше & \cite{Pouille2001} & виміряно \\
4 & Заборона $a\wedge\bar b$~\eqref{eq:veto} з АБО й сталою одиницею функціонально повна (твердження~\ref{prop:gabacomplete}) & Доведення в розділі. Класичний результат: мережа порогових елементів із гальмівними входами реалізує будь-який логічний вираз. Твердження про модель, досліду немає & \cite{McCullochPitts1943} & теорія \\
5 & Заборона із затримкою дає детектор напрямку руху з оптимальною швидкістю $v^*=\Delta x/d$~\eqref{eq:vstar} & Сітківка кроля: вибірковість до напрямку створюється гальмуванням, яке за руху в «нульовому» напрямку гасить збудження. Окремий запис збуджувального й гальмівного входів вибіркових клітин показав, що зірчасті амакринові клітини (\nt{GABA}) гальмують їх несиметрично --- лише відростками, спрямованими в «нульовий» бік. Формула $v^*$ --- висновок розділу & \cite{Barlow1965,Fried2002} & виміряно \\
6 & Шунтувальне гальмування ділить напругу, а не віднімає~\eqref{eq:shunt} (рис.~\ref{fig:shunt}) & Закон Ома для подільника з трьох провідностей за рівноважного потенціалу хлору біля спокою; для нейрона без імпульсів & висновок у розділі & теорія \\
7 & На частоту імпульсів шунт діє відніманням, а ділення виходить від гальмування разом зі врівноваженим фоновим шумом & Модель: у нейрона, що спрацьовує, струм через шунт майже сталий, і ефект на частоту віднімальний. Дослід: у пірамідні нейрони соматосенсорної кори щура (зрізи) через динамічну фіксацію вводили фонові збуджувальні й гальмівні провідності; одночасне врівноважене зростання обох ділить крутість відповіді без помітного зсуву частоти. Огляд: ділення на загальну активність трапляється в багатьох відділах мозку & \cite{HoltKoch1997,Chance2002,Carandini2012} & виміряно \\
8 & При $\beta>1$ взаємне гальмування обирає одного переможця (твердження~\ref{prop:wta}, \eqref{eq:wta}) & Власні значення $-(1\pm\beta)/\tau$ --- висновок у розділі. Частоти $0.73$ і $0.53$ при $\beta=0.5$ --- нерухома точка $r_{1,2}=(I_{1,2}-\beta I_{2,1})/(1-\beta^2)$; скрипту в \texttt{models/} немає. Прямого досліду в розділі не наведено & висновок у розділі & теорія \\
9 & Скидання пам'яті сигналом через \nt{GABA}; дві групи із взаємним гальмуванням --- засувка & Випливає з рис.~\ref{fig:bistable} і твердження~\ref{prop:wta}. Досліду немає & висновок у розділі & теорія \\
10 & Рівновага петлі \nt{GLUT}--\nt{GABA} стійка, якщо гальмування досить сильне й досить швидке (твердження~\ref{prop:ei}) & Модель двох популяцій~\eqref{eq:ei}; умови (i) і (ii) --- визначник і слід її матриці, виведені в розділі & \cite{WilsonCowan1972} & теорія \\
11 & При $w_{EE}=1.5$ і $w_{EI}w_{IE}=2$ швидке гальмування ($\tau_I=2\ms$) тримає $E^*=0.67h$, повільне ($30\ms$) розгойдує коливання (рис.~\ref{fig:ei}) & Чисельний розв'язок~\eqref{eq:ei}, скрипт \texttt{figures/make\_ei.py} & розрахунок & модель \\
12 & Кора працює в режимі стабілізації гальмуванням; підпис --- підштовхнеш \nt{GABA}-нейрони, і їхня частота впаде~\eqref{eq:isn} & Теорія передбачила парадоксальну відповідь. Дослід: внутрішньоклітинні записи в первинній зоровій корі кішки --- при пригніченні відповіді подразником навколо гальмівна провідність не зростає, а падає разом зі збуджувальною. Оптогенетичне збудження PV-нейронів миші в зоровій, соматосенсорній і моторній корі знижує частоту гальмівних нейронів, зокрема без подразника й за легкого наркозу. Коефіцієнт $-1/3$ за чисел рис.~\ref{fig:ei} --- висновок розділу & \cite{Tsodyks1997isn,Ozeki2009,Sanzeni2020} & виміряно \\
13 & Петля «\nt{GLUT} збуджує \nt{GABA}, \nt{GABA} гальмує \nt{GLUT}» породжує швидкий ритм із періодом $12$--$50\ms$ & Оптогенетична стимуляція швидких \nt{GABA}-нейронів у соматосенсорній корі миші на частотах $8$--$200$\,Гц вибірково посилює гамма-ритм ($20$--$80$\,Гц, період $12$--$50\ms$); стимуляція \nt{GLUT}-нейронів посилює лише повільніші ритми & \cite{Cardin2009,Buzsaki2012} & виміряно \\
\bottomrule
\end{longtable}}
